\subsection{Worked solutions on the ten soils}
The ten soils of the ensemble differ widely in their retention. This section works one question on each of them, so that the reader can see how the same formulas give different answers.

\exbox{Ten-soil example 1a: Coarse loam (S01)}{Coarse loam (S01): $\alpha=0.0961$ cm$^{-1}$, $n=1.353$, AWC $=1.5$ mm\,cm$^{-1}$. Find $m$, $S_e$ at field capacity ($h=100$ cm) and wilting point ($h=15\,849$ cm), the suction at half depletion, and TAW, RAW ($p=0.5$) and the days of supply at $ET_c=5$ mm\,d$^{-1}$ for a root zone of 0.6 m.}\sol $m=1-1/1.353=0.261$ \ans $S_e(100)=[1+(0.0961\times100)^{1.353}]^{-0.261}=0.445$; $S_e(15\,849)=0.0752$ \ans Half depletion is $f=0.5$, found by bisection of Eq.~(\ref{eq:fdep}): $h_{0.5}=471$ cm \ans $TAW=1.5\times600/10\ldots=90$ mm, i.e.\ $1.50\times60$ cm; $RAW=0.5\times90=45$ mm \ans Supply: $RAW/ET_c=45/5=9.0$ days \ans

\exbox{Ten-soil example 1a: Coarse sand (S02)}{Coarse sand (S02): $\alpha=0.0536$ cm$^{-1}$, $n=1.600$, AWC $=0.5$ mm\,cm$^{-1}$. Find $m$, $S_e$ at field capacity ($h=100$ cm) and wilting point ($h=15\,849$ cm), the suction at half depletion, and TAW, RAW ($p=0.5$) and the days of supply at $ET_c=5$ mm\,d$^{-1}$ for a root zone of 0.6 m.}\sol $m=1-1/1.600=0.375$ \ans $S_e(100)=[1+(0.0536\times100)^{1.600}]^{-0.375}=0.356$; $S_e(15\,849)=0.0175$ \ans Half depletion is $f=0.5$, found by bisection of Eq.~(\ref{eq:fdep}): $h_{0.5}=303$ cm \ans $TAW=0.5\times600/10\ldots=30$ mm, i.e.\ $0.50\times60$ cm; $RAW=0.5\times30=15$ mm \ans Supply: $RAW/ET_c=15/5=3.0$ days \ans

\exbox{Ten-soil example 1a: Sand (S03)}{Sand (S03): $\alpha=0.1450$ cm$^{-1}$, $n=2.680$, AWC $=0.6$ mm\,cm$^{-1}$. Find $m$, $S_e$ at field capacity ($h=100$ cm) and wilting point ($h=15\,849$ cm), the suction at half depletion, and TAW, RAW ($p=0.5$) and the days of supply at $ET_c=5$ mm\,d$^{-1}$ for a root zone of 0.6 m.}\sol $m=1-1/2.680=0.627$ \ans $S_e(100)=[1+(0.1450\times100)^{2.680}]^{-0.627}=0.011$; $S_e(15\,849)=0.0000$ \ans Half depletion is $f=0.5$, found by bisection of Eq.~(\ref{eq:fdep}): $h_{0.5}=151$ cm \ans $TAW=0.6\times600/10\ldots=36$ mm, i.e.\ $0.60\times60$ cm; $RAW=0.5\times36=18$ mm \ans Supply: $RAW/ET_c=18/5=3.6$ days \ans

\exbox{Ten-soil example 1a: Loamy sand (S04)}{Loamy sand (S04): $\alpha=0.1240$ cm$^{-1}$, $n=2.280$, AWC $=0.9$ mm\,cm$^{-1}$. Find $m$, $S_e$ at field capacity ($h=100$ cm) and wilting point ($h=15\,849$ cm), the suction at half depletion, and TAW, RAW ($p=0.5$) and the days of supply at $ET_c=5$ mm\,d$^{-1}$ for a root zone of 0.6 m.}\sol $m=1-1/2.280=0.561$ \ans $S_e(100)=[1+(0.1240\times100)^{2.280}]^{-0.561}=0.040$; $S_e(15\,849)=0.0001$ \ans Half depletion is $f=0.5$, found by bisection of Eq.~(\ref{eq:fdep}): $h_{0.5}=172$ cm \ans $TAW=0.9\times600/10\ldots=54$ mm, i.e.\ $0.90\times60$ cm; $RAW=0.5\times54=27$ mm \ans Supply: $RAW/ET_c=27/5=5.4$ days \ans

\exbox{Ten-soil example 1a: Sandy loam (S05)}{Sandy loam (S05): $\alpha=0.0750$ cm$^{-1}$, $n=1.890$, AWC $=1.2$ mm\,cm$^{-1}$. Find $m$, $S_e$ at field capacity ($h=100$ cm) and wilting point ($h=15\,849$ cm), the suction at half depletion, and TAW, RAW ($p=0.5$) and the days of supply at $ET_c=5$ mm\,d$^{-1}$ for a root zone of 0.6 m.}\sol $m=1-1/1.890=0.471$ \ans $S_e(100)=[1+(0.0750\times100)^{1.890}]^{-0.471}=0.165$; $S_e(15\,849)=0.0018$ \ans Half depletion is $f=0.5$, found by bisection of Eq.~(\ref{eq:fdep}): $h_{0.5}=217$ cm \ans $TAW=1.2\times600/10\ldots=72$ mm, i.e.\ $1.20\times60$ cm; $RAW=0.5\times72=36$ mm \ans Supply: $RAW/ET_c=36/5=7.2$ days \ans

\exbox{Ten-soil example 1a: Loam (S06)}{Loam (S06): $\alpha=0.0360$ cm$^{-1}$, $n=1.560$, AWC $=1.7$ mm\,cm$^{-1}$. Find $m$, $S_e$ at field capacity ($h=100$ cm) and wilting point ($h=15\,849$ cm), the suction at half depletion, and TAW, RAW ($p=0.5$) and the days of supply at $ET_c=5$ mm\,d$^{-1}$ for a root zone of 0.6 m.}\sol $m=1-1/1.560=0.359$ \ans $S_e(100)=[1+(0.0360\times100)^{1.560}]^{-0.359}=0.466$; $S_e(15\,849)=0.0286$ \ans Half depletion is $f=0.5$, found by bisection of Eq.~(\ref{eq:fdep}): $h_{0.5}=332$ cm \ans $TAW=1.7\times600/10\ldots=102$ mm, i.e.\ $1.70\times60$ cm; $RAW=0.5\times102=51$ mm \ans Supply: $RAW/ET_c=51/5=10.2$ days \ans

\exbox{Ten-soil example 1a: Silt loam (S07)}{Silt loam (S07): $\alpha=0.0200$ cm$^{-1}$, $n=1.410$, AWC $=2.0$ mm\,cm$^{-1}$. Find $m$, $S_e$ at field capacity ($h=100$ cm) and wilting point ($h=15\,849$ cm), the suction at half depletion, and TAW, RAW ($p=0.5$) and the days of supply at $ET_c=5$ mm\,d$^{-1}$ for a root zone of 0.6 m.}\sol $m=1-1/1.410=0.291$ \ans $S_e(100)=[1+(0.0200\times100)^{1.410}]^{-0.291}=0.686$; $S_e(15\,849)=0.0943$ \ans Half depletion is $f=0.5$, found by bisection of Eq.~(\ref{eq:fdep}): $h_{0.5}=483$ cm \ans $TAW=2.0\times600/10\ldots=120$ mm, i.e.\ $2.00\times60$ cm; $RAW=0.5\times120=60$ mm \ans Supply: $RAW/ET_c=60/5=12.0$ days \ans

\exbox{Ten-soil example 1a: Sandy clay loam (S08)}{Sandy clay loam (S08): $\alpha=0.0590$ cm$^{-1}$, $n=1.480$, AWC $=1.4$ mm\,cm$^{-1}$. Find $m$, $S_e$ at field capacity ($h=100$ cm) and wilting point ($h=15\,849$ cm), the suction at half depletion, and TAW, RAW ($p=0.5$) and the days of supply at $ET_c=5$ mm\,d$^{-1}$ for a root zone of 0.6 m.}\sol $m=1-1/1.480=0.324$ \ans $S_e(100)=[1+(0.0590\times100)^{1.480}]^{-0.324}=0.417$; $S_e(15\,849)=0.0375$ \ans Half depletion is $f=0.5$, found by bisection of Eq.~(\ref{eq:fdep}): $h_{0.5}=369$ cm \ans $TAW=1.4\times600/10\ldots=84$ mm, i.e.\ $1.40\times60$ cm; $RAW=0.5\times84=42$ mm \ans Supply: $RAW/ET_c=42/5=8.4$ days \ans

\exbox{Ten-soil example 1a: Clay loam (S09)}{Clay loam (S09): $\alpha=0.0190$ cm$^{-1}$, $n=1.310$, AWC $=1.6$ mm\,cm$^{-1}$. Find $m$, $S_e$ at field capacity ($h=100$ cm) and wilting point ($h=15\,849$ cm), the suction at half depletion, and TAW, RAW ($p=0.5$) and the days of supply at $ET_c=5$ mm\,d$^{-1}$ for a root zone of 0.6 m.}\sol $m=1-1/1.310=0.237$ \ans $S_e(100)=[1+(0.0190\times100)^{1.310}]^{-0.237}=0.753$; $S_e(15\,849)=0.1704$ \ans Half depletion is $f=0.5$, found by bisection of Eq.~(\ref{eq:fdep}): $h_{0.5}=618$ cm \ans $TAW=1.6\times600/10\ldots=96$ mm, i.e.\ $1.60\times60$ cm; $RAW=0.5\times96=48$ mm \ans Supply: $RAW/ET_c=48/5=9.6$ days \ans

\exbox{Ten-soil example 1a: Silty clay loam (S10)}{Silty clay loam (S10): $\alpha=0.0100$ cm$^{-1}$, $n=1.230$, AWC $=1.7$ mm\,cm$^{-1}$. Find $m$, $S_e$ at field capacity ($h=100$ cm) and wilting point ($h=15\,849$ cm), the suction at half depletion, and TAW, RAW ($p=0.5$) and the days of supply at $ET_c=5$ mm\,d$^{-1}$ for a root zone of 0.6 m.}\sol $m=1-1/1.230=0.187$ \ans $S_e(100)=[1+(0.0100\times100)^{1.230}]^{-0.187}=0.878$; $S_e(15\,849)=0.3118$ \ans Half depletion is $f=0.5$, found by bisection of Eq.~(\ref{eq:fdep}): $h_{0.5}=906$ cm \ans $TAW=1.7\times600/10\ldots=102$ mm, i.e.\ $1.70\times60$ cm; $RAW=0.5\times102=51$ mm \ans Supply: $RAW/ET_c=51/5=10.2$ days \ans

